\paragraph{Prior on $\theta$.}
The dispersion parameter $\theta$ is strictly positive. When it is estimated,
the implementation introduces an unconstrained raw variable $\widetilde u\in\mathbb R$
and defines
\[
u
=
u_{\max}\tanh\!\left(\frac{\widetilde u}{u_{\max}}\right),
\qquad
\theta=\exp(u),
\qquad u_{\max}>0.
\]
Thus $\theta$ is positive and smoothly bounded between
$\exp(-u_{\max})$ and $\exp(u_{\max})$. The configuration currently exposes
$u_{\max}$ under the historical name \texttt{u\_clip}, although no hard clipping
is performed.

A Gaussian prior is placed on the raw variable,
\[
\widetilde u\sim\mathcal N(\mu_{\widetilde u},\sigma_u^2).
\]
If the requested center in effective log-dispersion space is $\mu_u$---for example,
$\mu_u=\log(\theta^{(0)})$ for a baseline value $\theta^{(0)}$---the raw prior
center is chosen by the inverse transformation,
\[
\mu_{\widetilde u}
=
u_{\max}\operatorname{arctanh}\!\left(\frac{\mu_u}{u_{\max}}\right).
\]
This guarantees that transforming the raw prior center gives exactly the requested
effective center. It requires $|\mu_u|<u_{\max}$.

